\section{Problem setup, HIVE geometry and the clean-entry criterion}
\label{sec:geometry}

% Sources: src/hive.py, data/hive_geometry.json, docs/01_hive_geometry.md,
% src/feasibility.py (Geometry.make, evaluate_gate, grid_from_candidates).
% Every worked number in this section was recomputed by running python against the
% repository code on 2026-10-30 (hive.Hive.from_json, Hive.classify_path,
% hive.classify_detail, feasibility.integrate_batch, feasibility.evaluate_gate).

Before any physics of spin can be judged, we need a precise answer to a purely geometric
question: \emph{what does a ball path have to look like to count as ``in''?} This section sets up
the coordinate frames, describes the HIVE cross-section that the model uses, defines the
\emph{clean-entry} test exactly as the code applies it, and derives from the geometry alone how
steeply the ball must be falling when it arrives from the unfavourable side. Everything later in
the paper (drag, Magnus lift, the launcher) only changes \emph{which paths} the robot can make;
this section decides \emph{which paths score}.

\subsection{The question in two dimensions}
\label{sec:geometry-question}

The target is the upward-facing CELL of the pivoting HIVE at the centre of the field. The HIVE is
one rigid beam with a box-shaped CELL at each end. It rests on a damper stop, so one CELL points up
and outward and the other points down and outward. A shot from the side the up CELL faces is
the \emph{favourable} side. A shot from the other side is the \emph{unfavourable} side: the up CELL's
roof is between the robot and the mouth. Our question is whether backspin from a single powered
wheel can bend a POLLEN path enough to score from the unfavourable side.

The model works in the vertical \emph{tipping plane}, the plane in which the CELLS rotate. It
uses the cross-section of the HIVE at its centre plane. A pure backspin shot fired in that plane
stays in it, so a two-dimensional (2D) model is enough for centred shots. Off-centre shots and the
CELL side walls are three-dimensional effects that this section treats only as a stated
limitation (Section~\ref{sec:geometry-assumptions}).

\subsection{Frames}
\label{sec:geometry-frames}

\paragraph{HIVE frame.} The origin is on the tile surface directly below the HIVE pivot axis, at
the HIVE centre plane. The axis $x$ (in \unit{m}) is horizontal and lies in the tipping plane, with
$+x$ pointing toward the audience side of the field. The axis $z$ (in \unit{m}) points up, and $y$
(in \unit{m}) runs along the pivot axis. The pivot axis is at height
\begin{equation}
  z_p = 1.11633\ \text{\unit{m}}\text{\confirmed},
  \label{eq:geometry-pivot}
\end{equation}
which is $43.95$\,\unit{in}\confirmed{} above the tiles (Competition Manual 9.6.1 and the official
field CAD). A point in the tipping plane is written $\vect{p} = (x, z)$.

\paragraph{Cell-local frame.} The HIVE body is stored once, in a frame that rotates with it. The
coordinate $u$ (in \unit{m}) runs along the cell axis, from the pivot outward toward the mouth of
the CELL that is up; $u > 0$ is the up CELL and $u < 0$ the down CELL. The coordinate $w$ (in
\unit{m}) is perpendicular to $u$ and points toward the CELL roof. Both CELLS have their roofs on
the $+w$ side, so the body is mirror-symmetric about $u = 0$. The lateral coordinate $s$ (in
\unit{m}) runs across the CELL opening; $s = 0$ is the centre plane, the plane used in this paper.

\paragraph{From cell-local to HIVE frame.} Let $\tau$ (in rad, quoted in degrees) be the
\emph{tilt}, the angle of the cell axis above the horizontal. Let $\sigma \in \{+1, -1\}$ be the
pose sign: $\sigma = +1$ when the up CELL is on the $+x$ side and $\sigma = -1$ when it is on the
$-x$ side. The code (\code{Hive.from\_json}, inner function \code{world}) places a cell-local point
$(u, w)$ in the HIVE frame in two steps. First it rotates by $\tau$:
\begin{equation}
  \begin{pmatrix} x' \\ z' - z_p \end{pmatrix}
  =
  \begin{pmatrix} \cos\tau & -\sin\tau \\ \sin\tau & \cos\tau \end{pmatrix}
  \begin{pmatrix} u \\ w \end{pmatrix}.
  \label{eq:geometry-rotation}
\end{equation}
Then it mirrors the horizontal coordinate when the up CELL is on the $-x$ side, $x = \sigma x'$.
Combining the two steps gives
\begin{align}
  x &= \sigma\,(u\cos\tau - w\sin\tau), \label{eq:geometry-world-x}\\
  z &= z_p + u\sin\tau + w\cos\tau. \label{eq:geometry-world-z}
\end{align}
Equations~\eqref{eq:geometry-world-x}--\eqref{eq:geometry-world-z} are used whenever the tilt or the pivot
height differ from their nominal values. At the nominal pose the code reads pre-computed
HIVE-frame segments from \code{data/hive\_geometry.json}; these agree with
\eqref{eq:geometry-world-x}--\eqref{eq:geometry-world-z} to the five decimals stored.

\paragraph{Shot frame.} The flight solver always fires toward $+x$. A shooter standing on the
$x < 0$ side therefore uses the HIVE frame unchanged. A shooter on the $x > 0$ side fires toward
$-x$, so the code mirrors the \emph{geometry} ($x \to -x$, \code{Hive.mirrored}) instead of the
trajectory. In the shot frame the robot releases the ball at
\begin{equation}
  \vect{p}_r = (x_r, z_r) = (-d,\; h),
  \label{eq:geometry-release}
\end{equation}
where $d$ (in \unit{m}) is the horizontal distance from the release point to the pivot plane
$x = 0$, and $h$ (in \unit{m}) is the height of the ball \emph{centre} above the tiles at release.

\subsection{The cross-section and its key points}
\label{sec:geometry-section}

Table~\ref{tab:geometry-local} lists the cell-local coordinates that set the mouth. All of them come from
the official field CAD (STEP file, parsed by the project), at the centre plane $s = 0$.

\begin{table}[h]
  \centering
  \caption{Cell-local coordinates of the up CELL at $s = 0$ (from \code{body\_local} in
  \code{data/hive\_geometry.json}).}
  \label{tab:geometry-local}
  \begin{tabular}{@{}lll@{}}
    \toprule
    Symbol & Meaning & Value \\
    \midrule
    $u_M$ & mouth plane (outer rib face) & $0.54508$\,\unit{m}\confirmed \\
    $u_B$ & back wall (back skin) & $0.23952$\,\unit{m}\confirmed \\
    $w_F$ & floor (top of the bottom skin) & $-0.03734$\,\unit{m}\confirmed \\
    $w_C$ & ceiling (roof underside at the pentagon peak) & $0.32004$\,\unit{m}\confirmed \\
    $w_T$ & roof top (envelope of the rib tops) & $0.33325$\,\unit{m}\confirmed \\
    \bottomrule
  \end{tabular}
\end{table}

\paragraph{Mouth length.} The mouth runs along the line $u = u_M$ from the lower lip
$(u_M, w_F)$ to the upper lip $(u_M, w_C)$. Its length is
\begin{equation}
  L = w_C - w_F = 0.32004 - (-0.03734) = 0.35738\ \text{\unit{m}}\text{\derived}.
  \label{eq:geometry-mouth-length}
\end{equation}
The CELL depth is $u_M - u_B = 0.30556$\,\unit{m}\derived{} (manual: about $12$\,\unit{in}\confirmed).

\paragraph{Worked example: the lips in the HIVE frame.} Take $\sigma = +1$ and
$\tau = 30^\circ$\confirmed, so $\cos\tau = 0.86603$\derived{} and $\sin\tau = 0.5$\derived.
For the lower lip, $(u, w) = (0.54508, -0.03734)$, equations
\eqref{eq:geometry-world-x}--\eqref{eq:geometry-world-z} give
\begin{align*}
  x &= 0.54508 \times 0.86603 - (-0.03734)\times 0.5\\
    &= 0.47205 + 0.01867 = 0.49072\ \text{\unit{m}}\text{\derived},\\
  z &= 1.11633 + 0.54508 \times 0.5 + (-0.03734)\times 0.86603\\
    &= 1.11633 + 0.27254 - 0.03234 = 1.35653\ \text{\unit{m}}\text{\derived}.
\end{align*}
That is $53.41$\,\unit{in}\derived, inside the Field Acceptance Checklist value of
$53.5 \pm 0.5$\,\unit{in}\confirmed. Running the same formula in python for the other key points
gives Table~\ref{tab:geometry-keypoints}.

\begin{table}[h]
  \centering
  \caption{Key points of the up CELL, pose \code{pose\_pos}, $\tau = 30^\circ$ (python, \code{hive.py}).}
  \label{tab:geometry-keypoints}
  \begin{tabular}{@{}lrrl@{}}
    \toprule
    Point & $x$ (\unit{m}) & $z$ (\unit{m}) & Official check \\
    \midrule
    Lower lip $(u_M, w_F)$ & $0.49072$\derived & $1.35653$\derived & $53.5$\,\unit{in}\confirmed{} (checklist) \\
    Upper lip $(u_M, w_C)$ & $0.31203$\derived & $1.66603$\derived & $65.6$\,\unit{in}\confirmed{} (Fig.\ 9-10) \\
    Roof top at mouth $(u_M, w_T)$ & $0.30543$\derived & $1.67747$\derived & highest point of the HIVE \\
    Roof top at back $(u_B, w_T)$ & $0.04081$\derived & $1.52469$\derived & -- \\
    \bottomrule
  \end{tabular}
\end{table}

\paragraph{Obstacles.} The code treats the HIVE as a set of solid line segments in the $(x, z)$
plane. Segment $j$ runs from $\vect{a}_j$ to $\vect{b}_j$ (both in \unit{m}). Each pose has
$16$\derived{} CELL segments (8 per CELL: back wall, roof outer, upper lip, roof inner, floor,
lower lip, underside, tube top), $4$\derived{} pivot-plate segments, and $9$\derived{} static
A-frame segments. Four of the A-frame segments are the legs. They slant across the field, so at
$s = 0$ only their top exists; the file flags them \code{projected}, and the code drops them by
default (\code{include\_projected=False}). The default obstacle set therefore has
$25$\derived{} segments ($29$\derived{} with the legs). The A-frame does not rotate with the HIVE.

\subsection{Poses, the mouth normal and the two sides}
\label{sec:geometry-poses}

\paragraph{Poses.} The HIVE has two resting poses. In \code{pose\_pos} the up CELL is on the
$+x$ (audience) side, $\sigma = +1$. In \code{pose\_neg} it is on the $-x$ side, $\sigma = -1$.
\code{pose\_neg} is the exact mirror image ($x \to -x$) of \code{pose\_pos}. Every TIP flips the pose.

\paragraph{Outward normal of the mouth.} The mouth lies on $u = u_M$, so the direction pointing
out of the CELL through the mouth is the cell-local unit vector $(1, 0)$. Rotating it with
\eqref{eq:geometry-rotation} gives $(\cos\tau, \sin\tau)$, and the mirror multiplies the $x$ part by
$\sigma$:
\begin{equation}
  \vect{n} = (\sigma\cos\tau,\; \sin\tau).
  \label{eq:geometry-normal}
\end{equation}
This is exactly the code's \code{normal = np.array([sp * ct, st])}. At $\tau = 30^\circ$\confirmed{}
and $\sigma = +1$, $\vect{n} = (0.86603, 0.5)$\derived: the mouth faces outward and $30^\circ$\derived{}
above the horizontal. The mouth \emph{line} is perpendicular to $\vect{n}$, and in \code{pose\_pos}
its direction from lower to upper lip is $\vect{m} = (-\sin\tau, \cos\tau) = (-0.5, 0.86603)$\derived.

\paragraph{Which way is ``in''?} A ball with velocity $\vect{v} = (v_x, v_z)$ (in \unit{m/s})
is moving into the CELL through the mouth only if its velocity points against the outward normal:
\begin{equation}
  \vect{v}\cdot\vect{n} < 0.
  \label{eq:geometry-inward}
\end{equation}

\paragraph{Favourable and unfavourable sides.} The \emph{favourable} side is the side the up CELL
is on (\code{favorable\_side}: \code{"pos"} for \code{pose\_pos}, \code{"neg"} for
\code{pose\_neg}). A ball fired from there travels toward the pivot, which is roughly
$-\vect{n}$, so \eqref{eq:geometry-inward} is easy to satisfy. The \emph{unfavourable} side is the
other one: the ball travels away from the shooter in the same direction as $\vect{n}_x$, so it can
only satisfy \eqref{eq:geometry-inward} by falling steeply (Section~\ref{sec:geometry-steep}). In the code
(\code{feasibility.Geometry.make}) the unfavourable case is \code{pose\_pos} used as it is, with the
shooter at $x = -d$ firing toward $+x$; the favourable case is the same HIVE mirrored. Results for
\code{pose\_neg} are the exact mirror image and are not computed separately.

\subsection{Tilt and its bounds}
\label{sec:geometry-tilt}

The nominal tilt is $\tau_0 = 30^\circ$\confirmed: Fig.\ 9-10 of the manual labels it, and the CAD
base tube measures $30.0^\circ$\confirmed. Real fields differ, so we bound it using the one
tilt-sensitive number the field checklist controls, the lower-lip height.

\paragraph{Step 1: lip height as a function of tilt.} From \eqref{eq:geometry-world-z} with
$(u, w) = (u_M, w_F)$:
\begin{equation}
  z_{\mathrm{lip}}(\tau) = z_p + u_M\sin\tau + w_F\cos\tau.
  \label{eq:geometry-zlip}
\end{equation}

\paragraph{Step 2: sensitivity.} Differentiating \eqref{eq:geometry-zlip} with $z_p$ held fixed,
\begin{equation}
  \frac{\mathrm{d}z_{\mathrm{lip}}}{\mathrm{d}\tau} = u_M\cos\tau - w_F\sin\tau .
  \label{eq:geometry-dzdtau}
\end{equation}
Comparing with \eqref{eq:geometry-world-x} shows that this is exactly $x_{\mathrm{lip}}$ for
$\sigma = +1$: the lip moves up at a rate equal to its horizontal distance from the pivot, as
expected for a rotation about the pivot. At $\tau_0$ it is $0.49072$\,\unit{m/rad}\derived{} $=
19.32$\,\unit{in/rad}\derived{} $= 0.3372$\,\unit{in/deg}\derived.

\paragraph{Step 3: invert for the tolerance.} The checklist allows
$\Delta z_{\mathrm{lip}} = \pm 0.5$\,\unit{in}\confirmed. With the pivot at its nominal height,
\begin{equation}
  \Delta\tau \approx \frac{\Delta z_{\mathrm{lip}}}{\mathrm{d}z_{\mathrm{lip}}/\mathrm{d}\tau}
  = \frac{0.5}{0.3372} = 1.48^\circ\text{\derived} \;\to\; \tau \in [28.5^\circ, 31.5^\circ]\text{\derived}.
  \label{eq:geometry-tilt-bound}
\end{equation}
If the pivot height is also off by the manual's general $\pm 1$\,\unit{in}\confirmed{} tolerance,
the lip can move $1.5$\,\unit{in}\derived{} for the same checklist reading, giving
$1.5/0.3372 = 4.45^\circ$\derived, rounded to the wide bounds $[25.5^\circ, 34.5^\circ]$\derived{}
(\code{low\_wide}, \code{high\_wide}). The code re-poses the body at any tilt with
\eqref{eq:geometry-world-x}--\eqref{eq:geometry-world-z}; Table~\ref{tab:geometry-tilt} in
Section~\ref{sec:geometry-steep} shows how the entry requirement moves with $\tau$.

\subsection{Configuration space: inflating the obstacles by the ball radius}
\label{sec:geometry-cspace}

A real ball touches a wall when its \emph{surface} reaches the wall. It is easier to track one
point, the ball centre, and grow every obstacle by the ball radius instead. This is the
\emph{configuration-space} idea, and it is how \code{hive.py} works. The POLLEN radius is
\begin{equation}
  R = 0.0355\ \text{\unit{m}}\text{\confirmed}
  \label{eq:geometry-radius}
\end{equation}
(manual 9.8: diameter ``approximately'' $7.1$\,\unit{cm}\confirmed, so about $\pm 1$\,\unit{mm}\assumed{}
is uncertain).

\paragraph{Distance to one segment.} For a ball centre $\vect{p}$ and segment $j$, the code
(\code{seg\_point\_dist}) first finds the nearest point's position along the segment and clips it to
the segment:
\begin{equation}
  \lambda_j^\ast = \operatorname{clip}\!\left(
  \frac{(\vect{p}-\vect{a}_j)\cdot(\vect{b}_j-\vect{a}_j)}{\lVert\vect{b}_j-\vect{a}_j\rVert^2},\;0,\;1\right),
  \qquad
  \vect{q}_j = \vect{a}_j + \lambda_j^\ast(\vect{b}_j-\vect{a}_j).
  \label{eq:geometry-nearest}
\end{equation}
The clip is what gives the inflated segment its rounded end caps: beyond an end, the nearest
point is the end itself. The distance is $\delta_j(\vect{p}) = \lVert\vect{p} - \vect{q}_j\rVert$ (in \unit{m}).

\paragraph{Clearance.} The clearance (in \unit{m}) is the smallest gap between the ball surface and
any obstacle:
\begin{equation}
  c(\vect{p}) = \min_j \delta_j(\vect{p}) - R .
  \label{eq:geometry-clearance}
\end{equation}
The ball is in \emph{contact} when $c(\vect{p}) < 0$ (strictly; $c = 0$ is a graze and not a contact).

\paragraph{Along a path.} A path is a list of centre positions $\vect{p}_0, \vect{p}_1, \dots$ at
the solver's time steps. The exact classifier (\code{Hive.classify\_path}) tests every \emph{chord}
$\vect{p}_k\vect{p}_{k+1}$ against every obstacle with the exact segment-to-segment distance
(\code{seg\_seg\_dist}), so contact between samples is not missed. The sweep
(\code{evaluate\_gate}) uses a precomputed clearance grid instead (\code{CSpaceRaster}), with cells
of side $\Delta = 1$\,\unit{mm}\assumed. It looks up the cell containing each sample. Because the
distance function changes by at most $1$\,\unit{m} per metre moved, the lookup error is at most
$\Delta/\sqrt{2} = 0.71$\,\unit{mm}\derived.

\subsection{The mouth gate and the inside region}
\label{sec:geometry-gate}

\paragraph{Inside region.} The code calls a ball centre \emph{inside} the CELL when it lies in the
quadrilateral $\mathcal{I}$ whose cell-local corners are
\begin{equation}
  \mathcal{I}:\quad u \in [\,u_B + R,\; u_M\,], \qquad w \in [\,w_F + R,\; w_C - R\,].
  \label{eq:geometry-inside}
\end{equation}
In words: behind the mouth plane, and at least one radius away from the back wall, the floor and the
ceiling. It is mapped to the HIVE frame with \eqref{eq:geometry-world-x}--\eqref{eq:geometry-world-z}. Note
that the edge on the mouth plane is \emph{not} moved in by $R$: the ball counts as inside as soon as
its centre is behind the mouth plane.

\paragraph{Gate.} The \emph{gate} is the edge of $\mathcal{I}$ that lies on the mouth plane
(\code{Hive.gate}), ordered low end first:
\begin{equation}
  \vect{A} = \mathrm{world}(u_M,\; w_F + R), \qquad \vect{B} = \mathrm{world}(u_M,\; w_C - R),
  \label{eq:geometry-gate}
\end{equation}
where $\mathrm{world}(\cdot)$ is \eqref{eq:geometry-world-x}--\eqref{eq:geometry-world-z}. Its length is
\begin{equation}
  L_g = (w_C - R) - (w_F + R) = L - 2R = 0.35738 - 0.071 = 0.28638\ \text{\unit{m}}\text{\derived}.
  \label{eq:geometry-gate-length}
\end{equation}
For \code{pose\_pos} at $\tau = 30^\circ$\confirmed{} the code gives
$\vect{A} = (0.47298, 1.38728)$\,\unit{m}\derived{} and $\vect{B} = (0.32979, 1.63529)$\,\unit{m}\derived.
The gate is the window the ball \emph{centre} must pass through; the ball surface then just clears
the lips.

\subsection{The clean-entry criterion}
\label{sec:geometry-clean}

\paragraph{Definition.} A path is \emph{clean} when the ball centre crosses the gate
$\overline{\vect{A}\vect{B}}$ moving inward, \eqref{eq:geometry-inward}, and the clearance
\eqref{eq:geometry-clearance} stayed non-negative at every earlier point of the path. Contact
\emph{after} the crossing (for example with the floor) does not change the label; whether such a
ball stays in is the job of the contact model. The exact classifier returns one of three labels:
\code{clean}; \code{hit:<name>}, the first obstacle touched before entry; or \code{miss}, when the
path reaches neither (for example it falls short onto the field).

\paragraph{How the sweep finds gate crossings.} Without wind, drag and Magnus forces do not depend
on position. So a launch with speed $v$ (in \unit{m/s}), elevation $\theta$ (in degrees) and spin
gives one path \emph{shape} $\vect{\Pi}(t) = (\Pi_x(t), \Pi_z(t))$ (in \unit{m}, starting at
$(0,0)$ at time $t = 0$\,\unit{s}), and the release point only shifts it:
\begin{equation}
  \vect{p}(t) = \vect{p}_r + \vect{\Pi}(t) = \bigl(-d + \Pi_x(t),\; h + \Pi_z(t)\bigr).
  \label{eq:geometry-translate}
\end{equation}
\code{evaluate\_gate} uses this to avoid re-integrating for every release point. For a fixed $h$,
it keeps every sample $k$ of every path that satisfies two conditions:
\begin{align}
  A_z \le h + \Pi_z(t_k) &\le B_z, \label{eq:geometry-band}\\
  \dot{\vect{\Pi}}(t_k)\cdot\vect{n} &< 0. \label{eq:geometry-band-inward}
\end{align}
The first condition says the sample is at a height where the gate is; the second is
\eqref{eq:geometry-inward}. Because the gate is a straight, non-horizontal segment, a height fixes a
unique point on it. Let $s_g$ be the fraction along the gate (0 at $\vect{A}$, 1 at $\vect{B}$):
\begin{equation}
  s_g = \frac{h + \Pi_z(t_k) - A_z}{B_z - A_z}, \qquad
  x_g = A_x + s_g\,(B_x - A_x).
  \label{eq:geometry-gate-fraction}
\end{equation}
Requiring the shifted sample to sit exactly on that gate point, $-d + \Pi_x(t_k) = x_g$, gives the
one release distance at which this sample enters:
\begin{equation}
  d = \Pi_x(t_k) - x_g .
  \label{eq:geometry-d-from-gate}
\end{equation}
Candidates outside the swept range widened by $0.06$\,\unit{m}\assumed{} at each end,
$d \in [0.49, 1.81]$\,\unit{m}\derived{} (grid $0.55$ to $1.75$\,\unit{m}\assumed), are dropped. Each
remaining candidate is clean if the raster clearance is $\ge 0$ at every sample from the first one
whose shifted $x$ has reached the left edge of the raster box up to and including sample $k$
(samples outside the box count as free). (For a path whose $v_x$ reverses, the check
starts at the first sample.) Two consecutive clean samples $k, k+1$ of the same path bound a
continuous interval of clean release distances, and \code{grid\_from\_candidates} marks every grid
value of $d$ inside it as clean. The exact classifier reaches the same label in a different way:
the first sample inside $\mathcal{I}$ before any chord with negative clearance. The two agree
because the only way into $\mathcal{I}$ from outside without contact is through the gate.

% CODE/CODE: evaluate_gate reports the descent angle as atan2(-vz, vx) (signed vx, so it can
% exceed 90 deg when backspin reverses vx); Hive.classify_path reports atan2(-vz, |vx|).
% They agree whenever vx > 0 at entry.

\paragraph{Worked example: one shot through the gate.} To check the formulas by hand we use a
vacuum path (no air at all), so that $\vect{\Pi}(t)$ is a parabola. Air is added in later
sections. Take $v = 7.0$\,\unit{m/s}\assumed, $\theta = 81.0^\circ$\assumed, $h = 0.737$\,\unit{m}\assumed{}
and $d = 1.00$\,\unit{m}\assumed, with $g = 9.80665$\,\unit{m/s\textsuperscript{2}}\confirmed{}
(standard gravity, \code{flight.G\_STANDARD}). Then
$v_x = v\cos\theta = 1.0950$\,\unit{m/s}\derived{} and $v_{z0} = v\sin\theta = 6.9138$\,\unit{m/s}\derived.

\emph{Step 1: crossings of the mouth line.} The centre is on the mouth line when
$(\vect{p}(t) - \vect{A})\cdot\vect{n} = 0$. With $\vect{p}(t) = (-d + v_x t,\; h + v_{z0}t - \tfrac12 g t^2)$
this is a quadratic in $t$:
\begin{equation}
  \underbrace{\bigl(-\tfrac12 g\,n_z\bigr)}_{c_2} t^2
  + \underbrace{\bigl(n_x v_x + n_z v_{z0}\bigr)}_{c_1} t
  + \underbrace{\bigl(n_x(-d - A_x) + n_z(h - A_z)\bigr)}_{c_0} = 0 .
  \label{eq:geometry-quadratic}
\end{equation}
Its roots are $t_1 = 0.5057$\,\unit{s}\predicted{} and $t_2 = 1.2911$\,\unit{s}\predicted.

\emph{Step 2: test each crossing.} At $t_1$ the ball is still rising
($v_z = +1.954$\,\unit{m/s}\predicted), $\vect{v}\cdot\vect{n} = +1.926$\,\unit{m/s}\predicted{} (outward),
and $s_g = 6.42$\predicted{}, far outside $[0, 1]$: the ball passes the extended mouth line high above
the HIVE. It fails both \eqref{eq:geometry-band} and \eqref{eq:geometry-band-inward}. At $t_2$ the centre is at
$(0.4138, 1.4897)$\,\unit{m}\predicted, $s_g = 0.413$\predicted{} (inside the gate),
$v_z = -5.748$\,\unit{m/s}\predicted{} and $\vect{v}\cdot\vect{n} = -1.926$\,\unit{m/s}\predicted{}
(inward). This is the entry.

\emph{Step 3: no earlier contact.} \code{hive.classify\_detail} on the exact parabola returns
\code{clean} with a minimum clearance of $30.5$\,\unit{mm}\predicted{} before entry, reached as the
ball passes the upper lip. The raster sweep (\code{evaluate\_gate} with \code{grid\_from\_candidates}) gives
$28.3$\,\unit{mm}\predicted{} for the same cell. The difference is not raster error: the grid
cell at $d = 1.00$\,\unit{m}\assumed{} takes the smaller margin of the two bracketing gate samples, which
sit at $d = 1.0025$\,\unit{m}\predicted{} and $0.9981$\,\unit{m}\predicted{} (one solver step of $\Delta t = 2$\,\unit{ms}\assumed{}, so
$v_x\,\Delta t \approx 2.2$\,\unit{mm}\derived{} apart); the exact classifier at
$d = 1.0025$\,\unit{m}\predicted{} gives $28.0$\,\unit{mm}\predicted.
% CHECK: re-run 2026-10-30. Raster margins 28.34 / 32.31 mm at d = 1.0025 / 0.9981 m;
% exact classify_path 27.99 / 32.35 mm. The raster lookup error is <= 0.71 mm either way.

\emph{Step 4: the apex.} The highest point is at $t = v_{z0}/g = 0.705$\,\unit{s}\predicted{},
$z = h + v_{z0}^2/(2g) = 3.174$\,\unit{m}\predicted. There $v_z = 0$ but the ball is still moving
sideways at $v_x = 1.095$\,\unit{m/s}\predicted; it is not stationary.

\emph{Step 5: the neighbours.} Keeping $v = 7.0$\,\unit{m/s}\assumed{} and changing only $\theta$:
$78.0^\circ$\assumed{} is a \code{miss}, $80.5^\circ$\assumed{} is \code{hit:up\_floor}
(clearance on the first touching chord $-1.6$\,\unit{mm}\predicted), $81.5^\circ$\assumed{} is
\code{hit:up\_roof\_outer} ($-3.9$\,\unit{mm}\predicted) and $84.0^\circ$\assumed{} is
\code{hit:pivot\_plate\_top}. The overlap reported at first contact depends on the time step of
the path samples, so only its sign is meaningful.
% CHECK: -1.6 / -3.9 mm from flight.simulate (vacuum, dt_max = 1 ms) + classify_detail; the
% exact parabola sampled at 0.1 ms gives -0.3 / -0.1 mm, at 2 ms -5.5 / -3.9 mm. The earlier
% draft values (-0.4 / -2.3 mm) were not reproducible.
In the sweep (\code{evaluate\_gate} with \code{grid\_from\_candidates}), on a sub-grid
($v$ from $5.0$\,\unit{m/s}\assumed{} to $9.0$\,\unit{m/s}\assumed{} in $0.1$\,\unit{m/s}\assumed{} steps,
$\theta$ from $60^\circ$\assumed{} to $89.75^\circ$\assumed{} in $0.25^\circ$\assumed{} steps) at this
$(h, d)$, $82$\predicted{} launches reach the gate and $47$\predicted{} of them are clean; the other
$35$\predicted{} first touch \code{up\_roof\_outer} ($22$\predicted) or \code{up\_upper\_lip}
($13$\predicted). The clean launches form a thin diagonal band in the $(v, \theta)$ plane. This is a
statement about the model in vacuum, not about how often a real launcher would score.

\subsection{Steepness on the unfavourable side: the minimum descent angle}
\label{sec:geometry-steep}

We now derive, from the geometry alone, how steeply the ball must fall to enter from the
unfavourable side. Work in the unfavourable shot frame: \code{pose\_pos}, $\sigma = +1$, shooter
at $x = -d$, ball moving toward $+x$.

\paragraph{Descent angle.} Define the descent angle $\varphi$ (in degrees) of the velocity, measured
below the horizontal and positive when falling, as the code does:
\begin{equation}
  \varphi = \operatorname{atan2}(-v_z,\; v_x), \qquad \vect{v} = \lVert\vect{v}\rVert\,(\cos\varphi,\; -\sin\varphi).
  \label{eq:geometry-descent}
\end{equation}

\paragraph{Step 1: the kinematic condition.} Insert \eqref{eq:geometry-descent} and
\eqref{eq:geometry-normal} into \eqref{eq:geometry-inward}:
\begin{equation}
  \vect{v}\cdot\vect{n} = \lVert\vect{v}\rVert\,(\cos\varphi\cos\tau - \sin\varphi\sin\tau)
  = \lVert\vect{v}\rVert\cos(\varphi + \tau) < 0 .
  \label{eq:geometry-vdotn}
\end{equation}
For a ball still moving toward $+x$ ($-90^\circ < \varphi < 90^\circ$) this needs
$\varphi + \tau > 90^\circ$, so
\begin{equation}
  \varphi > \varphi_{\mathrm{kin}} = 90^\circ - \tau = 60^\circ\text{\derived}.
  \label{eq:geometry-phi-kin}
\end{equation}
This is the ``steeper than $60^\circ$'' rule. It only says the ball is moving into the plane of the
mouth; it says nothing yet about fitting past the lips. It is convenient to measure the excess
over this angle,
\begin{equation}
  \alpha = \varphi + \tau - 90^\circ ,
  \label{eq:geometry-alpha}
\end{equation}
so that \eqref{eq:geometry-phi-kin} reads $\alpha > 0$.

\paragraph{Step 2: the roof.} The highest point of the HIVE is the roof top at the mouth,
$\vect{T} = (0.30543, 1.67747)$\,\unit{m}\derived{} (Table~\ref{tab:geometry-keypoints}). Every gate point
has $x > 0.3298$\,\unit{m}\derived{} and the release has $x = -d < 0$, so the centre must cross the
vertical line $x = x_T$. At that moment the centre is directly above $\vect{T}$ (below it is inside
the HIVE), so the clearance \eqref{eq:geometry-clearance} is non-negative only if
\begin{equation}
  z(x_T) \ge z_T + R = 1.67747 + 0.0355 = 1.7130\ \text{\unit{m}}\text{\derived}.
  \label{eq:geometry-roof}
\end{equation}
So the path must rise at least $1.7130 - h$ above the release point. In vacuum this needs a vertical
launch speed of at least $\sqrt{2g(1.7130 - h)}$, which is $4.375$\,\unit{m/s}\derived{} at
$h = 0.737$\,\unit{m}\assumed. Drag raises this bound; backspin lift can lower it. Both are treated in
later sections.

\paragraph{Step 3: a straight line near the mouth.} Close to the mouth the path is almost straight.
In the worked example the ball reaches the gate at $5.85$\,\unit{m/s}\predicted{} and
$\varphi = 79.2^\circ$\predicted. Under gravity alone its curvature is
$\kappa = g\cos\varphi/\lVert\vect{v}\rVert^2 = 0.054$\,\unit{m\textsuperscript{-1}}\derived, a radius
of $18.7$\,\unit{m}\derived. Over a $0.2$\,\unit{m}\assumed{} chord near the mouth the path departs
from a straight line by about $0.3$\,\unit{mm}\derived. We therefore treat the last part of the path as
a straight line with descent angle $\varphi$.

\paragraph{Step 4: go to the cell-local frame.} Rotate the direction of motion into the cell-local
frame with the inverse of \eqref{eq:geometry-rotation}. Using $\varphi + \tau = 90^\circ + \alpha$,
\begin{equation}
  \hat{\vect{e}}_{uw}
  = \bigl(\cos(\varphi+\tau),\; -\sin(\varphi+\tau)\bigr)
  = (-\sin\alpha,\; -\cos\alpha).
  \label{eq:geometry-local-dir}
\end{equation}
In the CELL's own frame the ball is falling almost parallel to the mouth plane, from the roof side
toward the floor side, while moving inward at the small rate $\sin\alpha$. Let the centre cross the
mouth plane at $(u_M, w_c)$, with $w_F + R \le w_c \le w_C - R$ (on the gate). Before the crossing
the centre was at
\begin{equation}
  \bigl(u_M + \lambda\sin\alpha,\; w_c + \lambda\cos\alpha\bigr), \qquad \lambda \ge 0 ,
  \label{eq:geometry-backray}
\end{equation}
with $\lambda$ (in \unit{m}) the distance travelled back along the path. It is outside the mouth
($u > u_M$) the whole time.

\paragraph{Step 5: clear the upper lip.} The obstacle nearest this ray is the inner corner of the
upper lip, $\vect{C} = (u_M, w_C)$; the rest of the lip and the roof are further away from it. The
distance from $\vect{C}$ to the line through $(u_M, w_c)$ in direction \eqref{eq:geometry-local-dir} is
the length of the cross product of $\vect{C} - (u_M, w_c) = (0,\; w_C - w_c)$ with the unit
direction:
\begin{equation}
  \delta_C = \bigl|\,0\cdot(-\cos\alpha) - (w_C - w_c)(-\sin\alpha)\,\bigr| = (w_C - w_c)\sin\alpha .
  \label{eq:geometry-dist-upper}
\end{equation}
The nearest point lies at $\lambda = (w_C - w_c)\cos\alpha \ge 0$, so it is \emph{before} the
crossing and counts. The clean test needs $\delta_C \ge R$:
\begin{equation}
  w_c \le w_C - \frac{R}{\sin\alpha}.
  \label{eq:geometry-upper-cond}
\end{equation}

\paragraph{Step 6: the lower lip does not bind before entry.} For the lower-lip corner
$\vect{D} = (u_M, w_F)$ the same calculation gives the nearest point at
$\lambda = (w_F - w_c)\cos\alpha < 0$, which is \emph{after} the crossing. Before the crossing the
closest the ball gets to $\vect{D}$ is at the crossing itself, $w_c - w_F \ge R$, which the gate
already guarantees. A ball crossing near the low end may touch the lower lip or floor just after it
enters, but by the definition in Section~\ref{sec:geometry-clean} it is already in.

\paragraph{Step 7: the minimum descent angle.} A clean straight path exists when
\eqref{eq:geometry-upper-cond} leaves room above the gate's low end, $w_C - R/\sin\alpha \ge w_F + R$, that is
$\sin\alpha \ge R/(w_C - w_F - R) = R/(L - R)$. Hence
\begin{equation}
  \boxed{\;\varphi > \varphi_{\mathrm{gate}} = 90^\circ - \tau + \arcsin\!\frac{R}{L - R}\;}
  \label{eq:geometry-phi-gate}
\end{equation}
With $R/(L-R) = 0.0355/0.32188 = 0.1103$\derived, $\arcsin(0.1103) = 6.33^\circ$\derived{} and
$\varphi_{\mathrm{gate}} = 66.33^\circ$\derived{} at $\tau = 30^\circ$\confirmed.

\paragraph{A stricter version: passing both lips.} If we also ask that the ball touch nothing until it
is completely past the mouth plane, the line must also stay $R$ from $\vect{D}$:
$\delta_D = (w_c - w_F)\sin\alpha \ge R$. Combined with \eqref{eq:geometry-upper-cond} this needs
$L\sin\alpha \ge 2R$:
\begin{equation}
  \varphi > \varphi_{\mathrm{fit}} = 90^\circ - \tau + \arcsin\!\frac{2R}{L}
  = 60^\circ + 11.46^\circ = 71.46^\circ\text{\derived}.
  \label{eq:geometry-phi-fit}
\end{equation}
The quantity $L\sin\alpha$ is the width of the mouth as seen along the path. This is the
``about $71.5^\circ$'' figure in the project documents. The code's clean test is
\eqref{eq:geometry-phi-gate}, not \eqref{eq:geometry-phi-fit}.
% CODE/DOC: docs/01_hive_geometry.md section 4 and the "reason" fields of
% data/hive_geometry.json give > ~71.5 deg (60 + asin(2R/0.3574)) as the minimum descent and
% a vertical-drop window of 0.108 m (|x| 0.347-0.455 m). That is the both-lips criterion
% (eq:geometry-phi-fit). The code's clean test (centre crosses the gate before any contact)
% admits straight paths from 66.33 deg and a vertical-drop window of 0.125 m (x 0.348-0.473 m).
% docs/03_clean_entry.md section 2.1 attributes clean gate descents below 71.5 deg (down to
% 68.3 deg) to path curvature; the curvature offset over the mouth is < 1 mm, so the gate
% definition, not curvature, explains the difference. This paper follows the code.

\paragraph{Window width.} The set of clean straight paths at a given $\alpha$ is
$w_c \in [\,w_F + R,\; w_C - R/\sin\alpha\,]$. Measured perpendicular to the path, its width (in \unit{m}) is
\begin{equation}
  W_\perp(\alpha) = \Bigl(w_C - \frac{R}{\sin\alpha} - w_F - R\Bigr)\sin\alpha = (L - R)\sin\alpha - R .
  \label{eq:geometry-window}
\end{equation}
For a vertical drop, $\varphi = 90^\circ$, so $\alpha = \tau$ and the window is horizontal:
$W_\perp = 0.32188 \times 0.5 - 0.0355 = 0.1254$\,\unit{m}\derived. Under the stricter both-lips
rule the same drop leaves $L\sin\tau - 2R = 0.1077$\,\unit{m}\derived. At $\varphi = 70^\circ$\assumed{}
($\alpha = 10^\circ$) the clean window is only $W_\perp = 0.0204$\,\unit{m}\derived, about $0.29$\derived{}
ball diameters ($0.57$\derived{} radii).

\paragraph{Worked check against the code.} We fired families of straight lines at the gate with
\code{Hive.classify\_path} (\code{pose\_pos}, HIVE frame, $400$\assumed{} entry points along the gate per angle)
and bisected on $\varphi$ for the smallest angle with any clean line. Table~\ref{tab:geometry-tilt} compares
the result with \eqref{eq:geometry-phi-gate}. The perpendicular widths found at $\varphi = 70^\circ$,
$72^\circ$ and $90^\circ$ were $0.0202$\predicted, $0.0312$\predicted{} and $0.1249$\,\unit{m}\predicted,
against $0.0204$\derived, $0.0314$\derived{} and $0.1254$\,\unit{m}\derived{} from
\eqref{eq:geometry-window}; the small shortfall comes from the discrete entry points, spaced
$0.72$\,\unit{mm}\derived{} along the gate, so up to one spacing is lost at each end of the window.
% CHECK: re-run 2026-10-30, (n_clean - 1) x spacing x |sin(gate, path)|; draft had
% 0.0201 / 0.0310 / 0.1245 m and a "0.5 mm sampling" that does not match 400 points on 0.286 m.
The vertical drop is clean for $x \in [0.348, 0.4726]$\,\unit{m}\predicted, from one radius right of
the upper lip ($0.31203 + 0.0355 = 0.3475$\,\unit{m}\derived) to the gate's low end
($A_x = 0.47298$\,\unit{m}\derived).

\begin{table}[h]
  \centering
  \caption{Descent-angle limits on the unfavourable side versus tilt. $\varphi_{\mathrm{kin}}$:
  \eqref{eq:geometry-phi-kin}; $\varphi_{\mathrm{gate}}$: \eqref{eq:geometry-phi-gate} (the code's criterion);
  $\varphi_{\mathrm{fit}}$: \eqref{eq:geometry-phi-fit}; ``code'' is the bisection with
  \code{Hive.classify\_path}; $z_{\mathrm{lip}}$ from the re-posed \code{Hive}.}
  \label{tab:geometry-tilt}
  \begin{tabular}{@{}rrrrrr@{}}
    \toprule
    $\tau$ & $z_{\mathrm{lip}}$ (\unit{m}) & $\varphi_{\mathrm{kin}}$ & $\varphi_{\mathrm{gate}}$ & code & $\varphi_{\mathrm{fit}}$ \\
    \midrule
    $25.5^\circ$\derived & $1.3173$\derived & $64.5^\circ$\derived & $70.83^\circ$\derived & -- & $75.96^\circ$\derived \\
    $28.5^\circ$\derived & $1.3436$\derived & $61.5^\circ$\derived & $67.83^\circ$\derived & $67.85^\circ$\predicted & $72.96^\circ$\derived \\
    $30.0^\circ$\confirmed & $1.3565$\derived & $60.0^\circ$\derived & $66.33^\circ$\derived & $66.35^\circ$\predicted & $71.46^\circ$\derived \\
    $31.5^\circ$\derived & $1.3693$\derived & $58.5^\circ$\derived & $64.83^\circ$\derived & $64.85^\circ$\predicted & $69.96^\circ$\derived \\
    $34.5^\circ$\derived & $1.3943$\derived & $55.5^\circ$\derived & $61.83^\circ$\derived & -- & $66.96^\circ$\derived \\
    \bottomrule
  \end{tabular}
\end{table}

A steeper HIVE (larger $\tau$) turns the mouth more toward the sky and lowers every limit one degree
per degree of tilt. The $\arcsin$ term does not depend on tilt because the CELL itself does not change
shape. The worked shot of Section~\ref{sec:geometry-clean} entered at $\varphi = 79.21^\circ$\predicted,
above all three limits at $\tau = 30^\circ$\confirmed.

\subsection{Release point and the robot height limit}
\label{sec:geometry-release}

The release point \eqref{eq:geometry-release} has two free numbers. The distance $d$ is measured from the
release point to the pivot plane $x = 0$ (not to the mouth). The field's inner perimeter wall is at
$|x| = 1.8009$\,\unit{m}\confirmed, so $d$ cannot exceed that. The sweep grid uses
$d$ from $0.55$\,\unit{m}\assumed{} to $1.75$\,\unit{m}\assumed{} in $0.05$\,\unit{m}\assumed{} steps.

The height $h$ is limited by the robot size rule. Manual rule R105 caps the expanded robot at
$H_{\max} = 29$\,\unit{in}\confirmed{} $= 0.7366$\,\unit{m}\confirmed{} measured from the field
surface. The project assumes that the ball centre at release is no higher than this:
\begin{equation}
  h \le H_{\max} = 0.7366\ \text{\unit{m}}\text{\assumed},
  \label{eq:geometry-hmax}
\end{equation}
and the sweep uses $h \in \{0.20, 0.35, 0.50, 0.65, 0.737\}$\,\unit{m}\assumed.
% CODE/DOC: data/hive_geometry.json robot_limits.release_height_note says "ball centre at release
% <= 0.737 m (ball inside the sizing volume)". If the whole ball had to be inside the volume the
% centre limit would be 0.7366 - R = 0.701 m. The code (SweepGrid.h) uses h = 0.737 m as the
% centre height, so this paper follows the code and flags the G416 risk.
% CHECK: the grid value 0.737 m (feasibility.SweepGrid.h) is H_max = 0.7366 m rounded up by
% 0.4 mm, so the top grid row sits 0.4 mm above eq:geometry-hmax; negligible but noted.
If instead the whole ball had to stay inside the sizing volume, the centre limit would be
$H_{\max} - R = 0.701$\,\unit{m}\derived; the code uses the former. The lowest point of the HIVE (the
down CELL's lower lip) is at $0.778$\,\unit{m}\derived{}, above the release, so every shot must rise to
pass over the roof, condition \eqref{eq:geometry-roof}.

\subsection{Assumptions}
\label{sec:geometry-assumptions}

\begin{itemize}
  \item The model is 2D, in the HIVE centre plane $s = 0$, where the mouth is tallest. Off-centre the
        pentagonal roof is lower. For a lateral offset $|s|$, \code{hive.py} lowers the roof segments
        and the top of the mouth by
        \begin{equation}
          \Delta w(s) = (h_0 - h_w)\,\frac{\min(|s|,\, s_{\mathrm{half}})}{s_{\mathrm{half}}},
          \label{eq:geometry-roof-drop}
        \end{equation}
        with peak height $h_0 = 0.3556$\,\unit{m}\confirmed, side-wall height
        $h_w = 0.19329$\,\unit{m}\confirmed{} and half-width $s_{\mathrm{half}} = 0.254$\,\unit{m}\confirmed.
        This is a 2D proxy; the side walls themselves are not modelled.
  \item The obstacles are solid line segments with no thickness of their own. The roof top and the
        underside are envelopes of the ribs, accurate to about $11$\,\unit{mm}\derived. The ballast screw near
        each CELL opening and the three preloaded NECTAR are not modelled.
  \item The four A-frame legs are left out by default, because at $s = 0$ only their tops exist.
  \item The ball is a rigid sphere of radius $R$; its centre is the tracked point.
  \item ``Clean'' means the centre crosses the gate inward with no earlier contact (an assumed
        criterion). Contact after the crossing, and whether the ball then stays in, belongs to the contact
        model. A planned bank shot off the roof is outside this channel (and is a G417 risk).
  \item The tilt is a fixed $30^\circ$\confirmed{} rest angle, with $\pm 1.5^\circ$\derived{} and
        $\pm 4.5^\circ$\derived{} bounds; the HIVE does not move during the shot.
  \item The descent-angle limits \eqref{eq:geometry-phi-gate} and \eqref{eq:geometry-phi-fit} treat the last few
        centimetres of the path as straight.
  \item The release height limit \eqref{eq:geometry-hmax} applies to the ball centre.
  \item The worked shot is in vacuum and illustrates the method only; its numbers are not a prediction
        for POLLEN in air.
\end{itemize}
